\section{Reproductibilité : fichiers, nombres et temps machine}
\label{sec:reproducibility}

\subsection{Environnement enregistré et formats}

Le calcul étendu a été exécuté sur CPU sous Windows, dans l'environnement
Python local du projet, avec tableaux NumPy \texttt{float64}. Les versions
du tableau~\ref{tab:environment} proviennent du fichier
\path{environment.json}, et non de versions supposées disponibles au
moment de la lecture du rapport. Le paquet porte la version interne
$0.0.0$, qui n'est pas une publication.

\begin{table}[htbp]
 \centering
 \caption{Environnement du calcul étendu archivé ; versions enregistrées,
 sans installation ni nouvelle expérience pour la rédaction.}
 \label{tab:environment}
 \begin{tabular}{ll}
\toprule
Élément & Valeur enregistrée \\
\midrule
Python & 3.14.6 \\
Plateforme & Windows-11-10.0.26200-SP0 \\
numpy & 2.5.3 \\
matplotlib & 3.11.2 \\
pytest & 9.1.1 \\
pde-image-denoising & 0.0.0 \\
Calcul & CPU, float64 \\
\bottomrule
\end{tabular}

\end{table}

Une tentative initiale d'import de \path{skimage.metrics} avait chargé
une DLL SciPy refusée par Windows Application Control. La restriction
n'a pas été contournée. Les mesures finales utilisent donc les formules
NumPy documentées en section~\ref{sec:metrics} ; les paquets optionnels
\texttt{scikit-image} et SciPy, présents après cette tentative, ne sont
pas utilisés par le calcul. Cette différence d'implémentation est
consignée pour éviter une attribution trompeuse à une bibliothèque.

Chaque ligne de \path{metrics.csv} identifie image, bruit, graine,
méthode, conductance, $K$, pas, checkpoint, temps de diffusion, quatre
scores et durée cumulative. Des hashes identifient les observations et
sorties. \path{summary.csv} regroupe sans perdre la définition des
métriques individuelles. Les fichiers NPZ archivent les entrées et les
sorties nécessaires aux exemples ; les PDF/PNG, tables, configuration,
environnement, journal et manifeste complètent la chaîne de provenance.

\subsection{Intégrité et modes de vérification}

Le manifeste énumère les fichiers de résultats, leurs tailles et SHA-256,
et relie la configuration utilisée à son empreinte. Les hashes d'arrays
sont calculés sur une représentation canonique \texttt{float64} de
boutisme fixé, avec forme et type inclus. Le vérificateur contrôle la
structure des CSV, les clés de regroupement, les graines, les observations
communes, les métriques recalculées et toutes les moyennes et dispersions.
Les statistiques de durée sont recalculées depuis les durées enregistrées,
jamais comparées à un nouveau chronométrage.

Deux constats doivent être distingués. Le compte rendu local archivé
\path{results/extended/verification.json} atteste une reproduction
\emph{stricte} dans l'environnement Windows enregistré : égalité exacte
et aucune différence de hash recalculé. L'audit indépendant communiqué
pour cette phase atteste une reproduction \emph{numérique}
inter-environnements, avec \texttt{rtol=1e-10} et \texttt{atol=1e-12}.
Ce second constat est celui de l'audit fourni, et non une vérification
relancée pendant la rédaction.

Le mode numérique conserve des contrôles exacts de structure et
d'intégrité des fichiers archivés, mais accepte de faibles différences
des quantités recalculées sous la règle
$|a-b|\leq10^{-12}+10^{-10}|b|$. Il signale leurs éventuels hashes
différents sans annoncer une reproduction exacte. L'intégrité d'un
ZIP ou d'un manifeste et la proximité numérique entre environnements
sont donc deux garanties différentes. Le tableau~\ref{tab:verification}
résume les contrôles locaux réellement enregistrés ; les $165$ tests
réussis appartiennent à la validation antérieure, sans relance de la
suite scientifique pour cette rédaction.

\begin{table}[htbp]
 \centering
 \caption{Contrôles consignés dans la vérification stricte locale de
 phase~4 ; durée et horodatage ne sont pas assimilés à des quantités
 déterministes.}
 \label{tab:verification}
 \begin{tabular}{lll}
\toprule
Contrôle & Niveau & Erreur d'amplitude ou maximale \\
\midrule
Mode propre discret & $17\times19$, 7 pas & $3.60822483\times10^{-16}$ \\
Spatial & 16 & $5.20069556\times10^{-4}$ \\
Spatial & 32 & $1.30111858\times10^{-4}$ \\
Temporel & 8 & $1.43420373\times10^{-4}$ \\
Temporel & 16 & $7.15660884\times10^{-5}$ \\
\bottomrule
\end{tabular}

\end{table}

\subsection{Lire correctement les durées}

La durée enregistrée du pipeline expérimental est
$44{,}407302$~s dans son environnement d'origine. Elle inclut calculs,
métriques, écritures et figures ; elle n'est pas une mesure universelle
de coût. La somme des durées \emph{terminales} des $540$ trajectoires
est $16{,}5264453$~s. Le chronométrage du solveur inclut les copies
d'états aux checkpoints, mais pas la génération du bruit, les métriques,
les hashes ou les figures. Les durées aux checkpoints sont cumulatives :
les additionner sur tous les checkpoints compterait plusieurs fois le
même calcul. Enfin, $t=n\dt$ est un temps de diffusion dans une convention
numérique, sans relation d'identité avec ces secondes machine.
